\documentclass[11pt]{article}
\usepackage[margin=1in]{geometry}
\usepackage{amsmath,amssymb,amsthm,mathtools,microtype}
\usepackage[hidelinks]{hyperref}
\usepackage[T1]{fontenc}

\newtheorem{theorem}{Theorem}
\newtheorem{lemma}{Lemma}
\newcommand{\A}{\mathcal A}
\newcommand{\one}{\mathbf 1}
\newcommand{\Z}{Z(\A)}
\newcommand{\posinv}{\A_{+}^{-1}}

\title{Complete Classification for a Power--Congruence Norm Identity}
\author{Discovery packet for arXiv:2403.07341}
\date{13 August 2026}

\begin{document}
\maketitle

\begin{center}
\fbox{\parbox{0.92\linewidth}{\textbf{Status.}
Candidate full theorem, likely valid, pending human review. The problem is
present in the raw arXiv TeX only \emph{after} \texttt{\string\end\{document\}}
and is absent from the rendered paper. The theorem below therefore solves an
orphaned drafting remnant, not a question visible in the published PDF.}}
\end{center}

\section{Source problem and provenance}

Hatori and Oi's \emph{Non-linear characterization of Jordan
$*$-isomorphisms via maps on positive cones of $C^*$-algebras},
arXiv:2403.07341, studies positive cones of unital $C^*$-algebras
\cite{HatoriOi}. Its raw source contains the following final problem. For
$a\in\posinv$ and $p>0$, put
\[
 \phi_{a,p}(x)=\bigl(a^{p/2}x^pa^{p/2}\bigr)^{1/p},
 \qquad x\in\posinv.
\]
Determine when
\begin{equation}\label{eq:source}
 \|\phi_{a,p}(x)a^{-1}\|=\|x\| \qquad(x\in\posinv).
\end{equation}
The raw note gives two endpoint answers: every $a$ works for $p=2$, whereas
for $p=1$ the condition is equivalent to centrality of $a$.

The source's \texttt{\string\end\{document\}} occurs immediately before this
``Final remark.'' Consequently, the 25-page arXiv PDF ends with the
bibliography and does not display the problem. The local file
\path{source_problem_excerpt.txt} preserves the exact raw-source passage.
The file \path{figures/raw_source_problem.png} is a screenshot of that block;
there can be no honest crop of the problem from \texttt{source\_paper.pdf}.

\section{Proof intuition}

The norm identity looks nonlinear, but positive elements turn a norm bound
into an order bound. After squaring and changing variables, it says exactly
that raising to a power $q=2/p$ preserves \emph{and reflects} every comparison
with the ray through $c=a^p$:
\[
 z\le tc \quad\Longleftrightarrow\quad z^q\le t^q c^q.
\]
If $q\ne1$, one direction of this equivalence gives local monotonicity of a
strictly convex power function at a suitable inverse. Nagy's centrality
theorem then forces $c$, and hence $a$, to be central. The exceptional value
$q=1$ is precisely $p=2$, where the identity follows by taking the square of
the norm. This mechanism also explains why the two cases already recorded in
the raw source are the two branches of the complete answer.

\section{Formal result}

Throughout, $\A$ is a unital $C^*$-algebra, $\Z$ is its center, and
$\posinv$ is its cone of positive invertible elements.

\begin{lemma}[lower local power rigidity]\label{lem:rigidity}
Let $d\in\posinv$ and $r>1$. If
\[
 0<w\le d \quad\Longrightarrow\quad w^r\le d^r
 \qquad(w\in\posinv),
\]
then $d$ is central.
\end{lemma}

\begin{proof}
Put $b=d^{-1}$. If $v\in\posinv$ and $b\le v$, then inversion reverses the
order, so $v^{-1}\le d$. The hypothesis gives
$v^{-r}\le d^r$; another inversion gives
\[
 b^r=d^{-r}\le v^r.
\]
Thus the strictly convex increasing function $f(t)=t^r$ on $(0,\infty)$ is
locally monotone at $b$: $b\le v$ implies $f(b)\le f(v)$. By Nagy's
centrality theorem \cite[Theorem~1]{Nagy}, $b$ is central. Hence so is
$d=b^{-1}$.
\end{proof}

\begin{theorem}[complete classification]\label{thm:main}
Let $a\in\posinv$ and $p>0$. Then \eqref{eq:source} holds for every
$x\in\posinv$ if and only if
\[
 \boxed{\quad p=2\quad\text{or}\quad a\in\Z.\quad}
\]
\end{theorem}

\begin{proof}
First assume \eqref{eq:source}. Set
\[
 c=a^p,\qquad q=\frac{2}{p}.
\]
For arbitrary $y\in\posinv$, substitute $x=y^{1/p}$ and square the norm in
\eqref{eq:source}. Since $c^{1/2}=a^{p/2}$ and
$c^{-q/2}=a^{-1}$, this gives
\begin{equation}\label{eq:normreduction}
 \left\|c^{-q/2}\bigl(c^{1/2}yc^{1/2}\bigr)^q c^{-q/2}\right\|
 =\|y\|^q.
\end{equation}

We claim that for every $z\in\posinv$ and $t>0$,
\begin{equation}\label{eq:orderiff}
 z\le tc \quad\Longleftrightarrow\quad z^q\le t^q c^q.
\end{equation}
Indeed, write $z=c^{1/2}yc^{1/2}$. All elements whose norms occur below are
positive, so \eqref{eq:normreduction} yields the chain of equivalences
\begin{align*}
 z\le tc
 &\Longleftrightarrow y\le t\one
 \Longleftrightarrow \|y\|\le t\\
 &\Longleftrightarrow
 \left\|c^{-q/2}z^q c^{-q/2}\right\|\le t^q
 \Longleftrightarrow c^{-q/2}z^q c^{-q/2}\le t^q\one\\
 &\Longleftrightarrow z^q\le t^q c^q.
\end{align*}

Suppose first that $q>1$. Taking $t=1$ in \eqref{eq:orderiff} shows that
$z\le c$ implies $z^q\le c^q$. Lemma~\ref{lem:rigidity}, with $d=c$ and
$r=q$, makes $c$ central.

Now suppose that $0<q<1$. Put $r=1/q>1$ and $d=c^q$. If
$0<w\le d$, set $z=w^r$. Then
$z^q=w\le c^q$. The reverse implication in \eqref{eq:orderiff} gives
$z\le c$; equivalently,
\[
 w^r\le c=d^r.
\]
Lemma~\ref{lem:rigidity} makes $d$ central, and consequently
$c=d^r$ is central. In both nonexceptional cases $q\ne1$, functional
calculus gives $a=c^{1/p}\in\Z$. Since $q=1$ exactly when $p=2$, necessity
is proved.

For sufficiency, first let $a$ be central. It commutes with every positive
$x$, so
\[
 \phi_{a,p}(x)=(a^p x^p)^{1/p}=ax,
 \qquad \phi_{a,p}(x)a^{-1}=x,
\]
for every $p>0$. Finally, if $p=2$, no commutativity is needed: writing
$B=(ax^2a)^{1/2}$,
\[
 \|Ba^{-1}\|^2
 =\|(Ba^{-1})^*(Ba^{-1})\|
 =\|a^{-1}B^2a^{-1}\|
 =\|x^2\|=\|x\|^2.
\]
This is \eqref{eq:source} and completes the proof.
\end{proof}

\section{Verification and adversarial checks}

\begin{itemize}
\item The proof uses both directions of \eqref{eq:orderiff}. The reverse
direction is essential when $p>2$ (equivalently $q<1$).
\item All substitutions are bijections of $\posinv$: $x\mapsto x^p$ and
$y\mapsto c^{1/2}yc^{1/2}$. Thus no positive invertible test elements are
lost.
\item Lemma~\ref{lem:rigidity} invokes Nagy's theorem only for the strictly
convex power $t^r$, $r>1$, on the open interval $(0,\infty)$, exactly within
its hypotheses.
\item Random $2\times2$ tests with noncentral diagonal $a$ and noncommuting
$x$ fail away from $p=2$ and agree at $p=2$. The script
\texttt{verify\_examples.py} is a non-proof sanity check.
\end{itemize}

\section{Novelty, limitations, and review recommendation}

A bounded search on 13 August 2026 covered the local run indexes and parsed
arXiv corpus; exact formulas and phrases from the raw problem; the paper
title and authors; and close web searches for the power-congruence identity.
No later paper or explicit resolution of this exact classification was found.
Nagy's 2019 theorem is a decisive prior ingredient, but it does not state the
norm classification; the reduction \eqref{eq:orderiff} is the new synthesis.
This is not an exhaustive novelty claim.

The principal provenance limitation is unusually important: the problem is
not in the rendered source paper. A reviewer should (i) check the
norm-to-order equivalence, especially the $q<1$ branch; (ii) confirm the use
of Nagy's Theorem~1; and (iii) decide whether an unrendered arXiv drafting
remnant merits presentation as an open-problem solution. Human review is not
yet complete.

\begin{thebibliography}{9}
\bibitem{HatoriOi}
O. Hatori and S. Oi,
\emph{Non-linear characterization of Jordan $*$-isomorphisms via maps on
positive cones of $C^*$-algebras}, arXiv:2403.07341 (2024).

\bibitem{Nagy}
G. Nagy,
\emph{Characterizations of centrality in $C^*$-algebras via local
monotonicity and local additivity of functions},
Integral Equations Operator Theory \textbf{91} (2019), Paper No.~28,
\href{https://doi.org/10.1007/s00020-019-2526-2}{doi:10.1007/s00020-019-2526-2}.
\end{thebibliography}

\end{document}
